\Needspace{10\baselineskip}
\subsection{Density roots from a closed factor}
\label{subsec:closed-factor-roots}

The threshold formula of \cref{lem:split-product-normalization} identifies a complete subspace of density roots with uniformly vanishing tails even when the ambient lower mass is zero. Use the product data $X_0,Y,N,H,C,V_0,\mu,R$ of \cref{sec:zero-lower-mass-product-equations}, for an arbitrary closed K\"ahler factor $(Y,\kappa)$, and put $c_0=CV_0$ and $A=\max_XR$. The spaces $\mathcal D_H$ and $\mathcal D_H^0$ are now taken on the $N$-dimensional product with reference volume $\mathrm dV$.

\Needspace{13\baselineskip}
\begin{proposition}\label{prop:closed-factor-density-roots}\index{Density space!closed factor}\index{Uniform tails!closed factor}
For real or complex $\xi\in L^N(Y,\kappa^k)$, put $J\xi=R^{1/N}\pi_Y^*\xi$. Then
\[
 \|J\xi\|_{L^N(\mathrm dV)}=\|\xi\|_{L^N(\kappa^k)},\qquad
 \rho_H(J\xi)=c_0^{-1/N}\|\xi\|_{L^N(\kappa^k)}.
\]
The image of $J$ is a closed linear subspace of $\mathcal D_H^0$. If $g=|\xi|^N$, then
\[
 \rho_H\bigl(J\xi\mathbf1_{\{|J\xi|>K\}}\bigr)
 \leq\left(\frac1{c_0}\int_{\{g>K^N/A\}}g\kappa^k\right)^{1/N}
       \longrightarrow0
 \qquad(K\to\infty).
\]
For $\mathrm dV=\mu$, so that $R=1$, the exact tail formula is
\[
 \rho_H\bigl(\pi_Y^*\xi\mathbf1_{\{|\xi|>K\}}\bigr)
   =\left(\frac1{c_0}\int_{\{|\xi|>K\}}|\xi|^N\kappa^k\right)^{1/N}.
\]
\end{proposition}
\begin{proof}
Since $|J\xi|^N=R|\xi|^N$, integration against $\mathrm dV$ gives the first norm identity. For nonzero $\xi$, part (2) of \cref{lem:split-product-normalization}, applied to $g=|\xi|^N$, gives the root-functional identity; for $\xi=0$ it holds by definition. Applying the two identities to $\xi-\eta$ yields
\[
 \|J\xi-J\eta\|_{\mathcal D_H}
       =(1+c_0^{-1/N})\|\xi-\eta\|_{L^N(\kappa^k)}.
\]
Bounded approximations of $\xi$ in $L^N$ have bounded images because $R$ is bounded, so $J\xi\in\mathcal D_H^0$. Completeness of $L^N(Y,\kappa^k)$ and the displayed identity make the image closed.

On $\{|J\xi|>K\}$ one has $g>K^N/A$, and hence
\[
 |J\xi|\mathbf1_{\{|J\xi|>K\}}
    \leq\bigl(Rg\mathbf1_{\{g>K^N/A\}}\bigr)^{1/N}.
\]
Monotonicity of the root functional and its identity for factor roots give the estimate, with the value zero when the tail vanishes. Integrability of $g$ makes the bound tend to zero; the $L^N$ tails also vanish. When $R=1$, the amplitude cutoff itself is pulled back from $Y$, so the same norm identity gives equality.
\end{proof}

For $Y=\mathbb P^k$ and $\mathrm dV=\mu$, take $p_\alpha=\pi_Y^*(q_\alpha^{1/N})$ with $q_\alpha$ from \cref{lem:projective-sublogarithmic-density}. Then $p_\alpha=J(q_\alpha^{1/N})\in\mathcal D_H^0$, but $p_\alpha\notin L^r(\mathrm dV)$ for every $r>N$. Moreover,
\[
 \rho_H\bigl(p_\alpha\mathbf1_{\{p_\alpha>K\}}\bigr)
       \asymp(\log K)^{-k(1-\alpha)/N}\qquad(K\to+\infty).
\]
Indeed, in the radial variable $r$ of the proof of that lemma, $q_\alpha\asymp r^{-k}(1-\log r)^{-1-k(1-\alpha)}$ as $r\downarrow0$. With $L=1-\log r$, the density level set $\{q_\alpha>K^N\}$ lies between two sets $\{L>c\log K\}$ for fixed positive constants $c$, once $K$ is large. Since $q_\alpha\kappa^k$ has radial density comparable to $L^{-1-k(1-\alpha)}\,\mathrm dL$ near $o$, integrating over $\{L>c\log K\}$ and taking the $N$th root of the exact tail formula gives the claimed tail. Finally, $p_\alpha\notin L^r(\mathrm dV)$ because $\int_Xp_\alpha^r\,\mathrm d\mu=\int_Yq_\alpha^{r/N}\kappa^k$ and $q_\alpha\notin L^{r/N}(\kappa^k)$.

These roots enlarge the examples in $\mathcal D_H^0$ beyond the subcritical $L^r$ classes. They give neither $\mathcal D_H^0=\mathcal D_H$ nor a general criterion for an exact equation at zero lower mass.
